\begin{abstract}
A language model that learns tasks one after another through adapters forgets, because every new
adapter changes what the model computes on old inputs. FK-LoRA gives every task one frozen set of
adapters and sends each input to one of them. An answer can then change only if its route changes,
so retention becomes a property of the rule that sets the route, and can be proved or measured for
that rule. Our main rule reads no instruction. It is a recogniser over all tasks, built on frozen
states of the model itself; it does not depend on the order in which tasks arrive, and it blocks no
task. With it come an audit, which estimates and bounds, for each task and from its held-out
inputs, the share of its answers that have changed, and a handicap, which limits in advance what a
later task can take from an older one. A simple bound says what cannot be avoided: without
instructions, any method loses on two tasks together at least the share
of inputs that they have in common with different right answers. Benchmarks hide this, because
their inputs begin with the instruction of the task. When it is removed, the two published methods
that we ran change \AudSeedsGainloraChanged{} and \AudSeedsInfloraChanged{} of \AudSeedsCompared{}
answers (\FiveGainChanged{} and \FiveInfChanged{} with it). This loss was not forced: no input of
these benchmarks occurs in two tasks, and the recogniser gives \AudDevLLamOneEnd{}\% of the test
inputs of 15 benchmark tasks to their task (development data), \AudScaleLongThreeEnd{}\% on a second
benchmark and \AudFreshFourthLEnd{}\% on 30 unseen tasks. Trained without instructions on one
sequence of five of these tasks (three seeds), plain adapters behind the recogniser change
\AudSeedsRidgeChanged{} of \AudSeedsCompared{} answers and end at \AudSeedsRidgeAPOne{} Rouge-L,
against \AudSeedsGainloraAPOne{} and \AudSeedsInfloraAPOne{} for the published methods. In that
comparison, with a tenth of their training epochs, most of the gap is learning. With their
published number of epochs (one seed) the two methods learn the tasks nearly as well as plain adapters
with the same epochs and then lose them: they forget \AudFullGainloraFTInt{} and \AudFullInfloraFTInt{} points and end
at \AudFullGainloraAPOne{} and \AudFullInfloraAPOne{}, where plain adapters behind the recogniser
end at \AudFullOursRidgeAPOne{} with the same number of epochs (\AudOursRidgeAPOne{} with a tenth of them, same seed); with their instructions left in, the same two methods forget \AudFullDefOriginalGainloraFTOne{} and \AudFullDefOriginalInfloraFTOne{} points. When inputs do carry an instruction, a second rule gives an exact statement: each task
claims a zone around the frozen word-vectors of its instruction, and we prove that the answer to an
input inside a zone never changes. On SuperNI with T5-base no answer to an earlier task changed
(\FiveCertChanged{} of \FiveCertCompared{} on five tasks, three seeds; \FifteenCertChanged{} of
\FifteenCertCompared{} on fifteen), but reworded instructions leave the zones: on machine
paraphrases of \WidenTasksEval{} real tasks they certify \WidenRefTfivePLow{}\% to
\WidenRefTfivePHigh{}\% of the inputs for the right task, and five attempts to widen them missed
their bars. The recogniser does not read the instruction, so no rewording can change its routes.
With instructions (five tasks, three seeds, machine-written rewordings) it ends
at \AudUThreeHeavyRidgeAPOne{} under a heavy rewording where the zones end at \AudUThreeHeavyCertAPOne{}; with
the original instructions the zones end higher (\AudUThreeOriginalCertAPOne{} against
\AudUThreeOriginalRidgeAPOne{}) and change no answer where the recogniser changes
\AudUThreeOriginalRidgeChanged{} of \AudUThreeCompared{}. Inputs that occur in two tasks can at best be split between them without an
instruction, the bound in advance is loose, and nothing here is in the published setting.
Stand-alone code is released.
\end{abstract}
